\begin{theorem}[有限群 gauge 表示群胚的圆维指数律]\label{thm:spg-boundary-gauge-groupoid-circle-law}
设 \(\Gamma\) 为有限无向多图，连通分支记为 \(\Gamma_1,\dots,\Gamma_c\)，并令
\[
b_0(\Gamma):=c,
\qquad
b_1(\Gamma):=\rank_{\ZZ}H_1(\Gamma;\ZZ).
\]
取有限群 \(G\)。
对每个连通分支选取一个基点，并将其基本群识别为自由群 \(F_{b_1(\Gamma_i)}\)。
定义表示群胚
\[
\mathsf{Rep}(\Gamma,G)
:=
\left[
\prod_{i=1}^{c}\Hom\!\bigl(F_{b_1(\Gamma_i)},G\bigr)
\Big/
G^{c}
\right],
\]
其中 \(G^c\) 通过分支上的同步共轭作用于对象集。
则其群胚基数满足
\[
\bigl|\mathsf{Rep}(\Gamma,G)\bigr|_{\mathrm{grpd}}
=
|G|^{\,b_1(\Gamma)-b_0(\Gamma)}.
\]
\leanverified{paper\_spg\_boundary\_gauge\_groupoid\_circle\_law\_seeds}
\leanverified{paper\_spg\_boundary\_gauge\_groupoid\_circle\_law\_package}
\leanverified{paper\_spg\_boundary\_gauge\_groupoid\_circle\_law}
\end{theorem}
\begin{proof}
对每个连通分支 \(\Gamma_i\)，自由群 \(F_{b_1(\Gamma_i)}\) 的表示空间满足
\[
\Hom\!\bigl(F_{b_1(\Gamma_i)},G\bigr)\cong G^{\,b_1(\Gamma_i)},
\]
故对象集基数为
\[
\prod_{i=1}^{c}\bigl|\Hom(F_{b_1(\Gamma_i)},G)\bigr|
=
\prod_{i=1}^{c}|G|^{b_1(\Gamma_i)}
=
|G|^{b_1(\Gamma)}.
\]
另一方面，\(G^c\) 以分支为单位作同步共轭，其阶为 \(|G|^c=|G|^{b_0(\Gamma)}\)。

对任意有限群 \(H\) 作用于有限集 \(Y\) 的 action groupoid \([Y/H]\)，其群胚基数满足
\[
|[Y/H]|_{\mathrm{grpd}}
=
\sum_{[y]}\frac{1}{|\operatorname{Stab}_H(y)|}
=
\frac{|Y|}{|H|},
\]
其中最后一步由轨道分解恒等式
\(
\sum_{[y]}|H|/|\operatorname{Stab}_H(y)|=|Y|
\)
给出。
将 \(Y=\prod_i\Hom(F_{b_1(\Gamma_i)},G)\)、\(H=G^c\) 代入，即得
\[
\bigl|\mathsf{Rep}(\Gamma,G)\bigr|_{\mathrm{grpd}}
=
\frac{|G|^{b_1(\Gamma)}}{|G|^{b_0(\Gamma)}}
=
|G|^{\,b_1(\Gamma)-b_0(\Gamma)}.
\]
证毕。
\end{proof}

\begin{corollary}[边界 gauge 数据的圆维容量下界]\label{cor:spg-boundary-gauge-groupoid-capacity}
设 \(f:\Omega_n\twoheadrightarrow X\) 为超立方体粗粒化，\(B_f\) 为其边界邻接多图。
固定有限群 \(G\)，并以
\[
\mathsf{Code}_f(G):=X\times \mathsf{Rep}(B_f,G)
\]
作为 gauge 约化后的边界码本。
若其有效容量按群胚基数计量，则
\[
\bigl|\mathsf{Code}_f(G)\bigr|_{\mathrm{grpd}}
=
|X|\cdot |G|^{\,b_1(B_f)-b_0(B_f)}.
\]
因此，在该 gauge-规范化计数口径下，若全部微态 \(2^n\) 需要由粗粒标签与边界 gauge 数据共同承载，则必要条件为
\[
2^n\le |X|\cdot |G|^{\,b_1(B_f)-b_0(B_f)}.
\]
等价地，
\[
b_1(B_f)\ \ge\ b_0(B_f)\ +\ \frac{n\log 2-\log|X|}{\log|G|}.
\]
于是，在 gauge 约化语义下控制边界有效态数的主导参数是 \(b_1(B_f)\)，而非仅是边数 \(A(f)\) 本身。
\leanverified{paper\_spg\_boundary\_gauge\_capacity\_seeds}
\leanverified{paper\_spg\_boundary\_gauge\_capacity\_package}
\leanverified{paper\_spg\_boundary\_gauge\_groupoid\_capacity}
\end{corollary}
\begin{proof}
由定理 \ref{thm:spg-boundary-gauge-groupoid-circle-law} 直接取 \(\Gamma=B_f\) 得
\[
\bigl|\mathsf{Rep}(B_f,G)\bigr|_{\mathrm{grpd}}
=
|G|^{\,b_1(B_f)-b_0(B_f)}.
\]
再乘以粗粒标签集 \(X\) 的基数即可得到 \(|\mathsf{Code}_f(G)|_{\mathrm{grpd}}\) 的闭式。
若该规范化容量需要覆盖全部 \(2^n\) 个微态，则必有上述不等式；其对 \(b_1(B_f)\) 的重写即为最后一式。证毕。
\end{proof}

\endinput
